\subsection{Composition Series and the H\"older Program}

\begin{exercise}[3.4.1]
    Prove that if $G$ is an Abelian simple group then $G \cong \z_p$ for some prime $p$ (do not assume $G$ is a finite group).
\end{exercise}

\begin{sol}
    Since $G$ is Abelian, every subgroup of $G$ is normal. Moreover, simplicity of $G$ implies that its only subgroups are 1 and $G$. If $G$ is the trivial group, then $G \cong Z_1$. Suppose $G$ is non-trivial, and let $g \in G$ be non-identity. Then $\gen g \leq G$. Since $g \neq 1$ and $G$ is simple, then $\gen g = G$, hence $G$ is cyclic.

    Assume, by way of contradiction, that $\abs G$ is not prime. We have two cases:
    \begin{itemize}
        \item Suppose $\abs G$ is infinite. Then $G \cong \z$ has infinitely many proper non-trivial subgroups, contradicting the simplicity of $G$.
        \item Suppose $\abs G$ is finite and composite. Then there exists some divisor $d$ of $\abs G$ such that $1 < d < \abs G$. Since $G$ is cyclic, then there exists a unique subgroup of $G$ of order $d$, contradicting the simplicity of $G$.
    \end{itemize}
    Therefore, $\abs G$ is prime, and $G \cong Z_p$ for some prime $p$.
\end{sol}

\begin{exercise}[3.4.2]
    Exhibit all $3$ composition series for $Q_8$ and all $7$ composition series for $D_8$. List the composition factors in each case.
\end{exercise}

\begin{sol}
    The composition series for $Q_8$ are as follows:
    \begin{align*}
        1 & \nsub \gen{-1} \nsub \gen i \nsub Q_8 \\
        1 & \nsub \gen{-1} \nsub \gen j \nsub Q_8 \\
        1 & \nsub \gen{-1} \nsub \gen k \nsub Q_8
    \end{align*}
    The composition factors for each series are isomorphic to $Z_2$.

    The composition series for $D_8$ are as follows:
    \begin{align*}
        1 & \nsub \gen s \nsub \gen{s, r^2} \nsub D_8 \\
        1 & \nsub \gen{r^2s} \nsub \gen{s, r^2} \nsub D_8 \\
        1 & \nsub \gen{r^2} \nsub \gen{s, r^2} \nsub D_8 \\
        1 & \nsub \gen{r^2} \nsub \gen r \nsub D_8 \\
        1 & \nsub \gen{r^2} \nsub \gen{rs, r^2} \nsub D_8 \\
        1 & \nsub \gen{rs} \nsub \gen{rs, r^2} \nsub D_8 \\
        1 & \nsub \gen{r^3s} \nsub \gen{rs, r^2} \nsub D_8
    \end{align*}
    with each composition factor isomorphic to $Z_2$.
\end{sol}

\begin{exercise}[3.4.3]
    Find a composition series for the quasidihedral group of order $16$ (cf. \exref{2.5.11}). Deduce that $QD_{16}$ is solvable.
\end{exercise}

\begin{sol}
    A clear composition series is $1 \nsub \gen{\sigma^4} \nsub \gen{\sigma^2} \nsub \gen{\sigma} \nsub QD_{16}$, with each composition factor isomorphic to $Z_2$. Since all composition factors are Abelian, then $QD_{16}$ is solvable.
\end{sol}

\begin{exercise}[3.4.4]
    Use Cauchy's Theorem and induction to show that a finite Abelian group has a subgroup of order $n$ for each positive divisor $n$ of its order.
\end{exercise}

\begin{sol}
    Let $\abs G = m$. If $m = 1$, the result is trivial. We now proceed by induction on $m$.

    Suppose that every finite Abelian group of order strictly less than $m$ has a subgroup of order $n$ for each positive divisor $n$ of its order. Let $n \in \zp$ where $n \mid m$. If $n$ is prime, then there exists $g \in G$ where $\abs g = n$ by Cauchy's Theorem. Then $\abs g \leq G$ is a subgroup of order $n$, and the result is true. 
    
    If $n$ is not prime, then $n = kp$ for some prime divisor $p$ of $n$. Since $p \mid m$, there exists an element $h \in G$ where $\abs h = p$ by Cauchy's Theorem. Then $\gen h$ has order $p$, and $G/\gen h$ is a finite Abelian group of order $m/p < m$. Since $n$ divides $m$, then $k$ divides $m/p$. By the inductive hypothesis, there exists a subgroup $\bar H$ of $G/\gen h$ such that $\abs{\bar H} = k$. By the Lattice Isomorphism Theorem, there exists a subgroup $H$ of $G$ containing $\gen h$ such that $H/\gen h = \bar H$. Then $\abs H = \abs{\bar H} \abs{\gen h} = kp = n$, hence $H$ is a subgroup of $G$ of order $n$. By induction, the result holds for all finite Abelian groups.
\end{sol}

\begin{spex}[3.4.5]
    Prove that subgroups and quotient groups of a solvable group are solvable.
\end{spex}

\begin{sol}
    Let $G$ be a solvable group, and let $H \leq G$. Then there exists a chain of subgroups
    \[1 = G_0 \nsub G_1 \nsub \cdots \nsub G_n = G\]
    such that $G_{i + 1}/G_i$ is Abelian. Define the subgroups
    \[H_i = H \cap G_i \quad \text{for each $0 \leq i \leq n$}\]
    Then $H_0 = H \cap G_0 = 1$ and $H_n = H \cap G_n = H$. It is clear that $H_i \leq H_{i + 1}$ for each $0 \leq i \leq n - 1$. To show that $H_i \nsub H_{i + 1}$, let $g \in H_{i + 1}$ and $x \in H_i$. Then $g \in G_{i + 1}$ and $x \in G_i$. Since $G_i \nsub G_{i + 1}$, then $gxg\inv \in G_i$. Since $g, x \in H$, then $gxg\inv \in H$, hence $gxg\inv \in H_i$. It follows that $H_i \nsub H_{i + 1}$. By the Second Isomorphism Theorem, we have
    \[H_{i + 1}/H_i \cong H_{i + 1}G_i/G_i \leq G_{i + 1}/G_i\]
    Since $G_{i + 1}/G_i$ is Abelian, then so is $H_{i + 1}/H_i$. Therefore, $H$ is solvable with the chain
    \[1 = H_0 \nsub H_1 \nsub \cdots \nsub H_n = H\]
    Now let $N \nsub G$. Consider the subgroups
    \[N_i = G_iN \quad \text{for each $0 \leq i \leq n$}\]
    Since it is clear that $N_i \leq N_{i + 1}$, we need to show that $N_i \nsub N_{i + 1}$. Let $g \in G_{i + 1}, h \in G_i$, and $n_1, n_2 \in N$. Then
    \[(gn_1)(hn_2)(gn_1)\inv = g n_1 h n_2 n_1\inv g\inv = g(n_1 h n_2 n_1\inv)g \inv\]
    Note that $n_1 hn_2n_1\inv \in N_i$ since $h \in G_i \leq G$ and $N \nsub G$. Since $G_i \nsub G_{i + 1}$, then $g(n_1 h n_2 n_1\inv)g \inv \in G_i \subseteq N_i$. It follows that $N_i \nsub N_{i + 1}$.

    Lastly, by the Lattice Isomorphism Theorem, then $N_i/N \nsub N_{i + 1}/N$, and the Third Isomorphism Theorem concludes that
    \[(N_{i + 1}/N)/(N_i/N) \cong N_{i + 1}/N_i \cong G_{i + 1}/(G_iN \cap G_{i + 1}) = G_{i + 1}/G_i\]
    Since $G_{i + 1}/G_i$ is Abelian, then so is $(N_{i + 1}/N)/(N_i/N)$. Therefore, $G/N$ is solvable with the chain
    \[1 = N_0/N \nsub N_1/N \nsub \cdots \nsub N_n/N = G/N. \qh\]
\end{sol}

\begin{exercise}[3.4.6]
    Prove part (1) of the Jordan--H\"older Theorem by induction on $|G|$.
\end{exercise}

\begin{sol}
    Let $G$ be a finite group. If $G = 1$, then the result is trivial. Suppose now $G$ is non-trivial and that the result holds for all groups of order less than $\abs G$. If $G$ is simple, then the composition series is $1 \nsub G$. Suppose now $G$ is not simple, and let $N$ be a maximal, normal subgroup of $G$. Since $\abs N < \abs G$, we may conclude by the inductive hypothesis that there exists a composition series
    \[1 = N_0 \nsub N_1 \nsub \cdots \nsub N_k = N\]
    for some $k \in \zp$. Since $N \nsub G$ and $G/N$ is simple, then the chain
    \[1 = N_0 \nsub N_1 \nsub \cdots \nsub N_k = N \nsub G\]
    is a composition series for $G$. By induction, the result holds for all finite groups.
\end{sol}

\begin{exercise}[3.4.7]
    If $G$ is a finite group and $H \trianglelefteq G$, prove that there is a composition series of $G$, one of whose terms is $H$.
\end{exercise}

\begin{sol}
    If $H = G$, then the result follows trivially, since the last term in the composition series must be $H$. Now suppose $H$ is a proper, normal subgroup of $G$. If $\abs G = 1$, the result is trivial. Suppose $\abs G > 1$, and assume the result is true for all groups with order less than $\abs G$.

    Since $H$ is proper, then $\abs H < \abs G$, hence we have a composition series
    \[1 = H_0 \nsub H_1 \nsub \cdots \nsub H_m = H\]
    by the inductive hypothesis. Now consider the quotient group $G/H$. Since $\abs{G/H} < \abs G$, then by the inductive hypothesis, there exists a composition series
    \[1 = K_0/H \nsub K_1/H \nsub \cdots \nsub K_n/H = G/H\]
    for some $n \in \zp$. By the Lattice Isomorphism Theorem, we have $H = K_0 \nsub K_1 \nsub \cdots \nsub K_n = G$. Since each $K_{i + 1}/K_i \cong (K_{i + 1}/H)/(K_i/H)$ is simple, then the chain
    \[1 = H_0 \nsub H_1 \nsub \cdots \nsub H_m = H = K_0 \nsub K_1 \nsub \cdots \nsub K_n = G\]
    is a composition series for $G$ with one of its terms equal to $H$. By induction, the result holds for all finite groups.
\end{sol}

\begin{exercise}[3.4.8]
    Let $G$ be a \emph{finite} group. Prove that the following are equivalent:
    \begin{enumerate}
        \item[(i)] $G$ is solvable.
        \item[(ii)] $G$ has a chain of subgroups
        \[1 = H_0 \trianglelefteq H_1 \trianglelefteq H_2 \trianglelefteq \cdots \trianglelefteq H_s = G\]
        such that $H_{i+1}/H_i$ is cyclic, $0 \le i \le s-1$.
        \item[(iii)] All composition factors of $G$ are of prime order.
        \item[(iv)] $G$ has a chain of subgroups
        \[1 = N_0 \trianglelefteq N_1 \trianglelefteq N_2 \trianglelefteq \cdots \trianglelefteq N_t = G\]
        such that each $N_i$ is a normal subgroup of $G$ and $N_{i+1}/N_i$ is Abelian, $0 \le i \le t-1$.
    \end{enumerate}
    \noindent [For (iv), prove that a minimal non-trivial normal subgroup $M$ of $G$ is necessarily Abelian and then use induction. To see that $M$ is Abelian, let $N \trianglelefteq M$ be of prime index (by (iii)) and show that $x^{-1}y^{-1}xy \in N$ for all $x,y \in M$ (cf. \exref{3.1.40}). Apply the same argument to $gNg^{-1}$ to show that $x^{-1}y^{-1}xy$ lies in the intersection of all $G$-conjugates of $N$, and use the minimality of $M$ to conclude that $x^{-1}y^{-1}xy = 1$.]
\end{exercise}

\begin{sol}
    \textbf{(i) \rightimp (ii)}: Suppose $G$ is solvable, and consider the chain of subgroups
    \[1 = G_0 \nsub G_1 \nsub \cdots \nsub G_n = G\]
    where $G_{i + 1}/G_i$ is Abelian. We prove the following lemma: if $A$ is a finite Abelian group, then there exists a chain of subgroups
    \[1 = A_0 \nsub A_1 \nsub \cdots \nsub A_k = A\]
    such that $A_{j + 1}/A_j$ is cyclic. We prove this lemma by induction on $\abs A$. If $\abs A = 1$, the result is trivial. Suppose $\abs A > 1$ and that the result holds for all Abelian groups of order less than $\abs A$. Since $A$ is finite Abelian, there exists some non-identity $a \in A$ such that $\gen a \nsub A$. Then $\abs{\gen a} < \abs A$, hence by the inductive hypothesis, there exists a chain of subgroups
    \[1 = B_0 \nsub B_1 \nsub \cdots \nsub B_m = \gen a\]
    such that $B_{j + 1}/B_j$ is cyclic. Now consider the quotient group $A/\gen a$. Since $\abs{A/\gen a} < \abs A$, then by the inductive hypothesis, there exists a chain of subgroups
    \[1 = C_0/\gen a \nsub C_1/\gen a \nsub \cdots \nsub C_\ell/\gen a = A/\gen a\]
    for some $\ell \in \zp$. By the Lattice Isomorphism Theorem, we have $\gen a = C_0 \nsub C_1 \nsub \cdots \nsub C_\ell = A$. Since each $C_{i + 1}/C_i \cong (C_{i + 1}/\gen a)/(C_i/\gen a)$ is cyclic, then the chain
    \[1 = B_0 \nsub B_1 \nsub \cdots \nsub B_m = \gen a = C_0 \nsub C_1 \nsub \cdots \nsub C_\ell = A\]
    satisfies the lemma. By induction, the lemma holds for all finite Abelian groups. 

    Applying this lemma to each Abelian quotient $G_{i + 1}/G_i$ in the original chain, we obtain a chain of subgroups
    \[1 = H_0 \nsub H_1 \nsub H_2 \nsub \cdots \nsub H_s = G\]
    such that $H_{i + 1}/H_i$ is cyclic, hence (ii) holds.

    \textbf{(ii) \rightimp (iii)}: Suppose $G$ has a chain of subgroups
    \[1 = H_0 \nsub H_1 \nsub H_2 \nsub \cdots \nsub H_s = G\]
    such that $H_{i + 1}/H_i$ is cyclic. We prove that each composition factor of $G$ is of prime order. We prove this by induction on $\abs G$. If $\abs G = 1$, the result is trivial. Suppose $\abs G > 1$ and that the result holds for all groups of order less than $\abs G$. Since $H_{s - 1} \nsub G$ and $G/H_{s - 1}$ is cyclic, then $G/H_{s - 1} \cong Z_n$ for some $n \in \zp$. If $n$ is prime, then the composition factors of $G/H_{s - 1}$ are of prime order. Since $\abs{H_{s - 1}} < \abs G$, then by the inductive hypothesis, the composition factors of $H_{s - 1}$ are also of prime order, hence all composition factors of $G$ are of prime order. Suppose now $n$ is not prime, so that $n = km$ for some $k, m \in \zp$ where $1 < k, m < n$. Then $Z_n$ has a proper, non-trivial subgroup $\gen{a^m}$ of order $k$. Let $\pi : G \to G/H_{s - 1}$ be the natural projection, and let $K = \pi\inv(\gen{a^m})$. Then $H_{s - 1} \nsub K \nsub G$ so that by the Lattice Isomorphism Theorem, $K/H_{s - 1} \cong \gen{a^m}$ is a proper, non-trivial subgroup of $G/H_{s - 1}$, contradicting that $H_{s - 1}$ is maximal normal in $G$. It follows that $n$ must be prime, hence all composition factors of $G$ are of prime order. By induction, the result holds for all finite groups.

    \textbf{(iii) \rightimp (iv)}: Suppose all composition factors of $G$ are of prime order. We prove that $G$ has a chain of subgroups
    \[1 = N_0 \nsub N_1 \nsub N_2 \nsub \cdots \nsub N_t = G\]
    such that each $N_i$ is a normal subgroup of $G$ and $N_{i + 1}/N_i$ is Abelian. Let $M$ be a minimal, non-trivial normal subgroup of $G$, and suppose $N \nsub M$ have prime index. Then for any $x, y \in M$, we have $x^{-1}y^{-1}xy \in N$ since $M/N$ is Abelian. Now for any $g \in G$, consider $gNg\inv \nsub gMg\inv = M$. By the same argument, we have $x^{-1}y^{-1}xy \in gNg\inv$ for all $x, y \in M$. It follows that $x^{-1}y^{-1}xy$ lies in the intersection of all $G$-conjugates of $N$. Since $M$ is minimal normal in $G$, then this intersection is either 1 or $M$. If it were $M$, then $M$ would be Abelian, contradicting that $M/N$ is non-trivial. Hence, $x^{-1}y^{-1}xy = 1$ for all $x, y \in M$, so that $M$ is Abelian.

    \textbf{(iv) \rightimp (i)}: This implication is immediate from the definition of solvable groups.
\end{sol}

\begin{exercise}[3.4.9]
    Prove the following special case of part (2) of the Jordan--H\"older Theorem: assume the finite group $G$ has two composition series
    \[ 1 = N_0 \trianglelefteq N_1 \trianglelefteq \cdots \trianglelefteq N_r = G \quad\text{and}\quad 1 = M_0 \trianglelefteq M_1 \trianglelefteq M_2 = G. \]
    Show that $r = 2$ and that the list of composition factors is the same. [Use the Second Isomorphism Theorem.]
\end{exercise}

\begin{sol}
    Consider $H = N_{r - 1} \cap M_1$. Then $H \nsub M_1$. Since $M_1/1 \cong M_1$ is simple, then either $H = 1$ or $H = M_1$. If $H = M_1$, then $N_{r - 1} \nsub M_1$. By simplicity of $M_1$, it must be that $N_{r - 1} = M_1$ so that $r = 2$. Both composition series are then the same, so we now suppose $H = 1$.

    Noting that $M_1 \nsub G$ and $N_{r - 1} \nsub G$, we may use the Second Isomorphism Theorem to conclude that
    \[N_{r - 1}M_1/M_1 \cong N_{r - 1}/(N_{r - 1} \cap M_1) = N_{r - 1}/1 \cong N_{r - 1}.\] Observe now that $M_1 \leq N_{r - 1}M_1 \leq G$, hence $N_{r - 1}M_1/M_1 \nsub G/M_1$. By simplicity of $G/M_1$, it must be that $N_{r - 1}M_1 = 1$ or $N_{r - 1}M_1 = G$. Since $N_{r - 1}$ is not trivial, it must be that $N_{r - 1}M_1 = G$. Now, $N_{r - 1}$ is simple, hence $N_{r - 2} = 1 = N_0$, and we may conclude that $r = 2$. Moreover, the composition factors are isomorphic:
    \[G/N_{r - 1} \cong M_1/(N_{r - 1} \cap M_1) \cong M_1 \quad \text{and} \quad N_{r - 1}/1 \cong N_{r - 1} \cong G/M_1 \qh\]
\end{sol}

\begin{exercise}[3.4.10]
    Prove part (2) of the Jordan--H\"older Theorem by induction on $\min\set{r,s}$. [Apply the inductive hypothesis to $H = N_{r-1} \cap M_{s-1}$ and use the preceding exercises.]
\end{exercise}

\begin{sol}
    Suppose $G$ has two composition series
    \begin{equation}
        \label{eqclubsuit}
        \tag{$\clubsuit$}
        1 = N_0 \nsub N_1 \nsub \cdots \nsub N_r = G
    \end{equation}
    and
    \begin{equation}
        \label{eqspadesuit}
        \tag{$\spadesuit$}
        1 = M_0 \nsub M_1 \nsub \cdots \nsub M_s = G.
    \end{equation}
    We induct on $\min\set{r,s}$. By the previous exercise, we may assume that $\min\set{r, s} > 2$. Suppose the statement is true for $\min\set{r, s} < k$ for some $k > 2$, and let $\min\set{r, s} = k$. Let $H = N_{r - 1} \cap M_{s - 1}$. If $N_{r - 1} \leq M_{s - 1}$ is proper, then by simplicity of $M_{s - 1}$, it must be that $N_{r - 1} = 1$, contradicting that $N_{r - 1}$ is non-trivial. Hence, $N_{r - 1} \nsub M_{s - 1}$. By a similar argument, we have $M_{s - 1} \nsub N_{r - 1}$. If we have $M_{s - 1} = N_{r - 1}$, then by the previous exercise, the composition factors of both series are the same, so we now suppose $M_{s - 1} \neq N_{r - 1}$ and they are not contained 
    in one another. They are then proper normal subgroups of $M_{s - 1}N_{r - 1}$. If $M_{s - 1}N_{r - 1} \neq G$, then $M_{s - 1}N_{r - 1} \nsub G$ is proper, so by simplicity of $G/N_{r - 1}$, it must be that $M_{s - 1}N_{r - 1} = N_{r - 1}$, contradicting that $M_{s - 1} \nsub N_{r - 1}$. Hence, $M_{s - 1}N_{r - 1} = G$. We may use the Second Isomorphism Theorem to conclude that
    \begin{equation}
        \label{eqheartsuit}
        \tag{$\heartsuit$}
        M_{s - 1}/H \cong M_{s - 1}N_{r - 1}/N_{r - 1} = G/N_{r - 1}
    \end{equation}
    and
    \begin{equation}
        \label{eqdiamondsuit}
        \tag{$\diamondsuit$}
        N_{r - 1}/H \cong M_{s - 1}N_{r - 1}/M_{s - 1} = G/M_{s - 1}.
    \end{equation}
    which shows that both $M_{s - 1}/H$ and $N_{r - 1}/H$ are simple. Consider now the composition series for $H$:
    \begin{equation}
        \label{eqstar}
        \tag{$\star$}
        1 = H_0 \nsub H_1 \nsub \cdots \nsub H_t = H.
    \end{equation}
    It is clear that appending $N_{r - 1}$ to the series \eqref{eqstar} gives a composition series for $N_{r - 1}$. Similarly, we get a composition series for $M_{s - 1}$. Moreover, there are $r - 2$ factors for the series for $N_{r - 1}$ and $s - 2$ factors for the series for $M_{s - 1}$. Since $\min\set{r - 1, s - 1} = k - 1 < k$, then by the inductive hypothesis, both series have the same number of factors and the same composition factors up to isomorphism. Therefore, $r - 2 = s - 2$ so that $r = s$. In \eqref{eqclubsuit} and \eqref{eqspadesuit}, both series have $r - 2$ factors isomorphic to \eqref{eqstar} up to some order, and the last 2 factors in \eqref{eqheartsuit} and \eqref{eqdiamondsuit} are isomorphic to $G/N_{r - 1}$ and $G/M_{s - 1}$ respectively. The result then follows by induction.
\end{sol}

\begin{exercise}[3.4.11]
    Prove that if $H$ is a non-trivial normal subgroup of the solvable group $G$ then there is a non-trivial subgroup $A$ of $H$ with $A \trianglelefteq G$ and $A$ Abelian.
\end{exercise}

\begin{sol}
    Since $G$ is solvable, we have the sequence
    \[1 = G_0 \nsub G_1 \nsub \cdots \nsub G_n = G\]
    where $G_{i + 1}/G_i$ is Abelian. Let $H_i = H \cap G_i$ for each $0 \leq i \leq n$. Since $H$ is non-trivial, consider the set of indices $I = \set{i \in \zp \mid G_i \cap H \neq 1}$. Since $H_n = H$, then $I$ is nonempty. By the Well Ordering Property, there exist some minimal index $k$ such that $H_k = G_k \cap H \neq 1$ but $H_{k - 1} = G_{k - 1} \cap H = 1$. Note that $H_k \nsub G$ since $H \nsub G$ and $G_k \nsub G$. Moreover, by the Second Isomorphism Theorem, we have
    \[H_k \cong H_k/1 = H_k/(H_k \cap G_{k - 1}) \cong H_k G_{k - 1}/G_{k - 1} \leq G_k/G_{k - 1}\]
    Since $G_k/G_{k - 1}$ is Abelian, then so is $H_k$. Therefore, $A = H_k$ is a non-trivial, Abelian normal subgroup of $G$ contained in $H$.
\end{sol}

\begin{exercise}[3.4.12]
    Prove (without using the Feit--Thompson Theorem) that the following are equivalent:
    \begin{enumerate}
    \item[(i)] Every group of odd order is solvable.
    \item[(ii)] The only simple groups of odd order are those of prime order.
    \end{enumerate}
\end{exercise}

\begin{sol}
    \textbf{(i) \rightimp (ii)} Suppose every group of odd order is solvable. Let $G$ be a simple group of odd order. Since $G$ is solvable, then by the characterization of solvable groups, all composition factors of $G$ are of prime order. Since $G$ is simple, then its only composition series is $1 \nsub G$, hence $G$ is of prime order.

    \noindent \textbf{(ii) \rightimp (i)} Suppose the only simple groups of odd order are those of prime order. Let $G$ be a group of odd order. We want to show that $G$ is solvable. Consider a composition series of $G$:
    \[1 = G_0 \nsub G_1 \nsub \cdots \nsub G_n = G.\]
    Each composition factor $G_{i+1}/G_i$ is simple and of odd order, so by assumption, each is of prime order and hence Abelian. Therefore, $G$ is solvable.
\end{sol}